\chapter{Fazit und Ausblick}\label{chap:Fazit-und-Ausblick}


\section{Zusammenfassung}

Die Evaluation zeigt, dass die Klassifikation der Dry-Bean-Sorten auf dem reduzierten Feature-Set gut funktioniert. Alle vier Modelle erreichen hohe und sehr ähnliche Werte. Die verbleibenden Fehler konzentrieren sich auf morphologisch ähnliche Sorten.

Die XAI-Ergebnisse ergänzen diese Bewertung. Sie zeigen, dass die Modellentscheidungen vor allem auf plausiblen Größen- und Formmerkmalen beruhen. Damit beantworten die Ergebnisse die Forschungsfrage grundsätzlich positiv: Die Modelle können die Bohnensorten zuverlässig klassifizieren, und die Entscheidungen des erklärten Random Forest sind mit XAI-Methoden inhaltlich nachvollziehbar. Einschränkungen bleiben durch korrelierte Features, ausgewählte lokale Fälle und die enge Auswahl der verwendeten Merkmale bestehen.


\section{Fazit}


\section{Ausblick}
